% WHAT A \read MAKES OF A LINE.
%
% \read builds a macro out of a line of a file, and the line is tokenized
% the way a line of a file is -- so all the rules a line has, it has:
%
%   an EMPTY line is a \par, the same as an empty line in a file. So is a
%   line of nothing but spaces, whose trailing spaces are stripped away
%   before \endlinechar goes on.
%
%   an INVALID character (catcode 15) is reported and forgotten, and a line
%   that has gone wrong is still READ TO ITS END -- so an invalid character
%   further along the same line is reported too.
%
%   an \outer name may no more stand in a definition here than anywhere
%   else. What has been read so far is shown, the name itself BECOMES A
%   SPACE, and a `}' is put in behind it:
%
%     Runaway definition?
%     ->a
%     ! Forbidden control sequence found while scanning definition of \a.
%
%   That `}' is read like any other. Where the line has a `{' still open it
%   closes it and the line is read on; where it has not, the rest of the
%   line is given up and so is the \read -- `a\lz' gives `a ' and nothing
%   more.
%
%   BRACES ARE COUNTED ACROSS LINES: a line that leaves one open makes the
%   \read take another, and another, until they balance or the file runs
%   out. At the end of the file the reference says
%
%     ! File ended within \read.
%     This \read has unbalanced braces.
%
%   and then takes the empty line the end of a file stands for, which is one
%   more \par -- put on AFTER the fault, so what the fault shows has not got
%   it.
%
% HSTeX gave an empty line nothing, gave up a line the moment it went wrong,
% and ended the \read where the reference put a brace in and read on. trip
% reads its own tripos.tex back with \uppercase made \outer and `0' made
% invalid, which is every one of these at once.
%
% Found and left open: the reference's context for all of this names the
% line being read -- `<read 3> gamma 0' -- and the brace it put in --
% `<inserted text> }'. HSTeX reads the line beside the source stack rather
% than on it, so neither frame is there to show.
\catcode`\{=1 \catcode`\}=2
\tracingonline=1
\lccode`\<=`\{ \lowercase{\gdef\lb{<}}
\lccode`\>=`\} \lowercase{\gdef\rb{>}} \def\space{ }
\immediate\openout9=zzread.dat
\immediate\write9{}
\immediate\write9{alpha \noexpand\zz 0 beta}
\immediate\write9{gamma \rb\space 0 delta}
\immediate\write9{plain 0 line}
\immediate\write9{\lb b\noexpand\zz c\lb d}
\immediate\write9{e}
\immediate\write9{   }
\immediate\closeout9
\outer\def\zz{}
\catcode`0=15
\openin3=zzread.dat
\message{[A an empty line]}\read3 to \a \show\a
\message{[B an outer name at balance nought, then an invalid character]}
\read3 to \a \show\a
\message{[C a closing brace too many, then an invalid character]}
\read3 to \a \show\a
\message{[D an invalid character on its own]}\read3 to \a \show\a
\message{[E an outer name inside braces, which reads on to the end]}
\read3 to \a \show\a
\closein3
\message{[done]}
\end
